% © 2026 Ing. David Jaroš — CC BY-NC-ND 4.0
%
% This work is licensed under a Creative Commons Attribution-NonCommercial-NoDerivatives
% 4.0 International License (CC BY-NC-ND 4.0).
%
% License History: Earlier drafts (up to v0.3) were released under CC BY 4.0.
% From v0.4 onward, all material is released under CC BY-NC-ND 4.0 to protect
% the integrity of the theoretical work during ongoing academic development.

% UBT-AI-PROVENANCE-BEGIN
% schema: ubt-ai-provenance/v1
% tier: B_machine_verified
% ai_assistance: disclosed
% human_review: machine-verification
% editorial_responsibility: Ing. David Jaroš
% policy: ../../AI_PROVENANCE.md
% notice: Machine-verified against named sources or verifiers; individual attestation is not claimed.
% UBT-AI-PROVENANCE-END

\documentclass[12pt,a4paper]{article}
\usepackage{amsmath,amssymb,amsthm}
\usepackage{mathtools}
\usepackage{geometry}
\usepackage{booktabs}
\usepackage{hyperref}
\usepackage{xcolor}
\usepackage{tcolorbox}
\tcbuselibrary{skins,breakable}
\geometry{left=2.5cm,right=2.5cm,top=2.5cm,bottom=2.5cm}

\newtheorem{theorem}{Theorem}[section]
\newtheorem{lemma}[theorem]{Lemma}
\newtheorem{proposition}[theorem]{Proposition}
\theoremstyle{remark}
\newtheorem{remark}[theorem]{Remark}

\newcommand{\Neff}{N_{\mathrm{eff}}}
\newcommand{\GSM}{G_{\mathrm{SM}}}

\newtcolorbox{statusbox}[1][]{colback=gray!8!white,colframe=black!60,
  arc=3pt,boxrule=1pt,fonttitle=\bfseries,title=Status Summary,#1}

\title{$N_{\mathrm{eff}}$ Reconciliation: Two Quantities or One?\\
       \large Resolving the Tension between Route~R2 and the Loop Audit}
\author{Ing.\ David Jaro\v{s}}
\date{2026-05-10}

\begin{document}
\maketitle

\begin{abstract}
Two independent analyses of UBT give superficially conflicting values of
$N_{\mathrm{eff}}$: Route~R2 (\texttt{neff12\_derivations.tex}) gives
$\Neff(\mathrm{R2})=\dim G_{\mathrm{SM}}=12$, while the loop audit
(\texttt{step2\_AUDIT.tex}) gives $\Neff(\mathrm{loop})=3$.
This document resolves the tension by showing that these two quantities
measure \emph{different physical objects}.
$\Neff(\mathrm{R2})=12$ is an algebraic count of gauge generators in
$\mathbb{C}\otimes\mathbb{H}$; $\Neff(\mathrm{loop})=3$ is the number
of charged complex scalar fields contributing to the one-loop
$U(1)_{\mathrm{EM}}$ vacuum polarisation from the minimal scalar kinetic
action $S_{\mathrm{kin}}[\Theta]$.  They are not in contradiction, but they
also cannot be equated without an additional physical mechanism.
\end{abstract}

\tableofcontents

%──────────────────────────────────────────────────────────────────────────────
\section{The Central Question}
\label{sec:question}
%──────────────────────────────────────────────────────────────────────────────

Route R2 in \texttt{research\_tracks/T3\_ALPHA/neff12\_derivations.tex}
(Theorem~2.2) states:
\begin{equation}
  \Neff(\mathrm{R2}) \;=\; \dim G_{\mathrm{SM}}
  \;=\; \underbrace{8}_{SU(3)_c}
       +\underbrace{3}_{SU(2)_L}
       +\underbrace{1}_{U(1)_Y}
  \;=\; 12.
  \label{eq:R2}
\end{equation}
The loop audit in \texttt{canonical/n\_eff/step2\_AUDIT.tex} concludes:
\begin{equation}
  \Neff(\mathrm{loop}) \;=\; N_{\mathrm{phases}} \;=\; 3
  \quad\Rightarrow\quad B_0 \;=\; 2\pi.
  \label{eq:audit}
\end{equation}
Are these two $\Neff$ values the same quantity measured differently,
or different quantities altogether?

%──────────────────────────────────────────────────────────────────────────────
\section{Step 1 — Precise Definitions}
\label{sec:definitions}
%──────────────────────────────────────────────────────────────────────────────

\subsection{$\Neff(\mathrm{R2})$: algebraic generator count}

\begin{definition}[$\Neff(\mathrm{R2})$]
$\Neff(\mathrm{R2}) \;=\; \dim G_{\mathrm{SM}}$ is the total number of
independent generators of the Standard Model gauge group
$G_{\mathrm{SM}} = SU(3)_c\times SU(2)_L\times U(1)_Y$ that are embedded
in $\mathcal{B} = \mathbb{C}\otimes\mathbb{H}$.
This is a pure algebraic fact, independent of any loop calculation or
field-equation dynamics.
\end{definition}

\subsection{$\Neff(\mathrm{loop})$: one-loop photon polarisation count}

\begin{definition}[$\Neff(\mathrm{loop})$]
$\Neff(\mathrm{loop})$ is the number of independent charged complex scalar
fields $\Phi_k$ that contribute to the one-loop photon vacuum polarisation
tensor $\Pi^{\mu\nu}(k)$ with a positive coefficient of
$e^2 \ln(\Lambda^2/m^2)$ in the Euclidean cut-off scheme, starting from
the minimal kinetic action
\begin{equation}
  S_{\mathrm{kin}}[\Theta]
  \;=\; \int d^4x\,d\psi\;
  \mathrm{Sc}\!\bigl[(D_\mu\Theta)^\dagger(D^\mu\Theta)\bigr].
  \label{eq:S_kin}
\end{equation}
\end{definition}

From the mode decomposition in \texttt{step1\_mode\_decomposition.tex}
(Theorem~3.1), after integrating over $\psi$, the action reduces to
\begin{equation}
  S_{\mathrm{eff}} \;=\; \sum_{k=1}^{3}
  \int d^4x\;\bigl|({\partial}_\mu + ieA_\mu)\Phi_k\bigr|^2,
  \qquad \Phi_k = \Theta_k,\; k=1,2,3.
  \label{eq:S_eff}
\end{equation}
Therefore $\Neff(\mathrm{loop}) = 3$.

%──────────────────────────────────────────────────────────────────────────────
\section{Step 2 — Do All SM Generators Contribute to U(1)$_{\mathrm{EM}}$
         Vacuum Polarisation?}
\label{sec:contributions}
%──────────────────────────────────────────────────────────────────────────────

The core assumption in Route~R2 is: "each generator corresponds to one
independent gauge-charged mode contributing to the $U(1)_{\mathrm{EM}}$
vacuum polarisation."  This step is examined critically below.

\subsection{SU(3)$_c$ sector (8 generators)}

Gluons are electrically neutral; they carry $\mathrm{SU}(3)_c$ colour but
zero $U(1)_{\mathrm{EM}}$ charge.  In standard QED, gluons do not appear
in the one-loop photon self-energy at all (they do in QCD corrections to
$\alpha_s$, but not to the photon propagator at one loop).

At the compactification scale where UBT lives, the SU(3) gauge sector
is associated with the $\mathbb{Z}_2\times\mathbb{Z}_2\times\mathbb{Z}_2$
involution subgroup of $\mathcal{B}$ (see
\texttt{canonical/su3\_derivation/su3\_from\_involutions.tex}).  These
involutions act on $\mathrm{Im}\,\mathbb{H}$ and are EM-neutral under the
diagonal $U(1)_{\mathrm{EM}}$ phase.

\begin{proposition}[SU(3) generators do not contribute to photon loop]
\label{prop:SU3}
None of the 8 SU(3)$_c$ generators embedded in $\mathcal{B}$ contribute
to the one-loop $U(1)_{\mathrm{EM}}$ vacuum polarisation from
$S_{\mathrm{kin}}[\Theta]$.
\end{proposition}

\begin{proof}
The SU(3) sector acts on the colour-triplet sub-space of the internal
$\mathbb{Z}_2^3$ structure.  In the minimal $U(1)_{\mathrm{EM}}$
covariant derivative $D_\mu = \partial_\mu + ieA_\mu$, the gauge field
$A_\mu$ couples only to the overall phase of $\Theta$.  The SU(3)
generators commute with this phase and hence with the minimal coupling.
Their virtual loop contributions to $\Pi^{\mu\nu}_{\mathrm{EM}}$ are zero
at leading order.
\end{proof}

\subsection{SU(2)$_L$ sector (3 generators)}

The $W^{\pm}$ bosons are EM-charged with $Q = \pm 1$.  However, they are
\emph{spin-1 vector fields}, not scalars.  In standard QFT, a charged
massive vector boson contributes to the photon $\beta$-function with the
opposite sign to a scalar:
\begin{equation}
  \Delta B_0^{W^\pm} = -\frac{11}{6}\,\frac{2\pi}{3}
  \quad\text{per massive }W^\pm\text{ (vector boson)}.
  \label{eq:W_contribution}
\end{equation}
The $Z^0$ boson is EM-neutral and does not contribute.

In the UBT action $S_{\mathrm{kin}}[\Theta]$ as written, the field $\Theta$
is treated as a complex scalar (not a vector field) coupling to
$U(1)_{\mathrm{EM}}$ via minimal coupling.  The $SU(2)_L$ gauge bosons
from $\mathbb{C}\otimes\mathbb{H}$ are \emph{not} included as independent
propagating fields in $S_{\mathrm{kin}}$ in the one-loop photon
vacuum polarisation computation.

\begin{remark}
If $W^\pm$ degrees of freedom were included as propagating charged scalars
(rather than vectors), they would contribute $+2 \times (2\pi/3)$ to $B_0$.
But their correct vector-boson contribution is $-2 \times (11/6)(2\pi/3)$,
which is negative and would \emph{decrease} $B_0$.  Neither scenario is
implemented in the audited scalar $S_{\mathrm{kin}}[\Theta]$.
\end{remark}

\subsection{U(1)$_Y$ sector (1 generator)}

The $B$-boson (hypercharge gauge boson) is EM-neutral at the level of
$U(1)_Y$ (before electroweak symmetry breaking, $U(1)_Y \neq U(1)_{\mathrm{EM}}$).
It does not contribute to the photon self-energy as an independent
propagating loop particle in the minimal coupling scheme.

%──────────────────────────────────────────────────────────────────────────────
\section{Step 3 — The Central Lemma}
\label{sec:lemma}
%──────────────────────────────────────────────────────────────────────────────

\begin{lemma}[Two different $\Neff$, not one]
\label{lem:two_Neff}
$\Neff(\mathrm{R2}) = 12$ and $\Neff(\mathrm{loop}) = 3$ are not the same
quantity:

\begin{enumerate}
  \item $\Neff(\mathrm{R2}) = 12$ counts \emph{algebraic generators of
        $G_{\mathrm{SM}}$ embedded in $\mathcal{B}$}.  It is an
        \textbf{[L0]} algebraic identity.  It does not require any
        assumption about which fields propagate in the photon loop.
  \item $\Neff(\mathrm{loop}) = 3$ counts \emph{charged complex scalar
        fields from $S_{\mathrm{kin}}[\Theta]$} that contribute to the
        one-loop $U(1)_{\mathrm{EM}}$ photon self-energy with the standard
        scalar-QED coefficient $e^2/(6\pi^2)$.  It is derived from
        \texttt{step1\_mode\_decomposition.tex} Theorem~3.1.
\end{enumerate}

The identification $\Neff(\mathrm{R2}) = \Neff(\mathrm{loop})$ requires
a separate mechanism showing that \emph{all} 12 generators contribute to
the photon loop via charged scalar-type excitations.  This mechanism is
\textbf{not proved} from $S_{\mathrm{kin}}[\Theta]$ in the current canonical
files.
\end{lemma}

\begin{proof}
From Section~\ref{sec:contributions}: SU(3) generators are EM-neutral
(Proposition~\ref{prop:SU3}); $W^\pm$ are vector bosons (negative
contribution); $Z^0$ and $B^0$ are EM-neutral.  No loop in
$S_{\mathrm{kin}}[\Theta]$ produces $\Neff = 12$ scalars.  The minimal
scalar action yields only the three $\Phi_k$ modes, giving
$\Neff(\mathrm{loop}) = 3$.  Therefore the two $\Neff$ values are not
equal.
\end{proof}

%──────────────────────────────────────────────────────────────────────────────
\section{Step 4 — Which $\Neff$ Enters $B_0$?}
\label{sec:which}
%──────────────────────────────────────────────────────────────────────────────

The quantity $B_0$ in
\begin{equation}
  \frac{1}{\alpha(\mu)}
  = \frac{1}{\alpha(\mu_0)}
  + \frac{B_0}{2\pi}\ln\frac{\mu}{\mu_0}
  \label{eq:running}
\end{equation}
is determined by which fields actually propagate in the one-loop photon
self-energy.

From $S_{\mathrm{kin}}[\Theta]$ alone:
\begin{equation}
  \boxed{
  \Neff(\mathrm{loop}) = 3
  \quad\Rightarrow\quad
  B_0 = \frac{2\pi}{3}\times 3 = 2\pi \approx 6.28.
  }
  \label{eq:B0_audit}
\end{equation}

If a future derivation shows that the full SM particle content
(quarks, leptons, $W^\pm$, etc.) emerges from $S[\Theta]$ with
the correct statistics and charges, then the effective $\Neff$ entering
$B_0$ at the EW scale would be the standard running-coupling answer
from the SM spectrum.  But that is a future problem, not the current one.

\begin{remark}[Impact on $n^*$]
With $B_0 = 2\pi$ and the stationarity equation $B = 2n_*/({\ln n_*+1})$:
\begin{equation}
  n_*(B_0 = 2\pi) \approx 10.54,
\end{equation}
which is not 137.  The alpha track T3\_ALPHA therefore remains
\textbf{CONDITIONAL-WEAK} until the mechanism connecting the full SM
particle content to $B$ is derived.
\end{remark}

%──────────────────────────────────────────────────────────────────────────────
\section{Consistency Check: Route R2 Is Internally Correct}
\label{sec:R2_correct}
%──────────────────────────────────────────────────────────────────────────────

Route R2 (\texttt{neff12\_derivations.tex} §R2) is \emph{internally} correct
as an algebraic statement: $\dim G_{\mathrm{SM}} = 12$ is a group-theory fact,
and the embedding of $G_{\mathrm{SM}}$ in $\mathcal{B}$ is proved at [L0].

The error is in the \emph{physical application}: the claim that each of the
12 generators "contributes to the $U(1)_{\mathrm{EM}}$ vacuum polarisation via
one independent gauge-charged mode" conflates the algebraic count with the
loop count.  These are different counting problems.

The coincidence $8 + 3 + 1 = 12 = 3\times 2\times 2$ is a genuine algebraic
coincidence explained in Remark~2.1 of \texttt{neff12\_derivations.tex}:
it follows from $\dim\,\mathrm{SU}(3) + \dim\,\mathrm{SU}(2) + \dim\,\mathrm{U}(1)
= N_{\mathrm{phases}}\times(N_{\mathrm{helicity}}\times N_{\mathrm{charge}})
= 3\times 4 = 12$.
This is a structural consistency, not a physical mechanism.

%──────────────────────────────────────────────────────────────────────────────
\section{Summary and Impact on WHAT\_IS\_PROVED.md}
\label{sec:summary}
%──────────────────────────────────────────────────────────────────────────────

\begin{center}
\begin{tabular}{lllll}
\toprule
Quantity & Value & What it counts & Proof level & Status \\
\midrule
$\Neff(\mathrm{R2})$ & 12 & $\dim G_{\mathrm{SM}}$ in $\mathcal{B}$ & [L0] & ✓ proved \\
$\Neff(\mathrm{loop})$ & 3 & Charged $\Phi_k$ from $S_{\mathrm{kin}}[\Theta]$ & [L1] & ✓ proved \\
$B_0$ from loop & $2\pi$ & Running from $\Neff(\mathrm{loop})$ & [L1] & ✓ proved \\
$\Neff(\mathrm{R2}) = \Neff(\mathrm{loop})$? & NOT EQUAL & — & [MC] OPEN & 
  identification unproved \\
\bottomrule
\end{tabular}
\end{center}

\medskip
\noindent
Updates to \texttt{WHAT\_IS\_PROVED.md}:
\begin{itemize}
  \item Row $\alpha$1b: $N_{\mathrm{eff}} = 12 = 3\times 2\times 2$:
    status remains \textbf{OPEN / [MC]} (loop identification unproved).
  \item Row $\alpha$1: $N_{\mathrm{phases}} = 3$ from $\mathrm{Im}\,\mathbb{H}$:
    status \textbf{[L0]} ✓ (not affected).
  \item New row: $\Neff(\mathrm{loop}) = 3$ from $S_{\mathrm{kin}}[\Theta]$:
    status \textbf{[L1]} ✓.
  \item New row: $B_0^{\mathrm{loop}} = 2\pi$: status \textbf{[L1]} ✓.
\end{itemize}

\begin{statusbox}
Výsledek: Two distinct $\Neff$ quantities — they are NOT the same.
$\Neff(\mathrm{R2}) = 12$ [L0] (algebraic); $\Neff(\mathrm{loop}) = 3$ [L1]
(from $S_{\mathrm{kin}}[\Theta]$).\\
Klasifikace: $\Neff(\mathrm{R2})$: [L0]; $\Neff(\mathrm{loop})$: [L1];
identification: [MC] OPEN.\\
Dopad na alpha track: $B_0 = 2\pi$ from the minimal scalar loop; $n^*\approx 10.54$.
T3\_ALPHA remains CONDITIONAL-WEAK until an extended loop derivation
produces $\Neff > 3$ from the full field content of $\mathcal{B}$.
\end{statusbox}

\end{document}
